\chapter{The language model as a reasoning component}
\section{From text to conditional prediction}
A language model operates on tokens: units produced by a tokenizer, which need not correspond to complete words. A name, identifier, or punctuation sequence may occupy several tokens. Tokenization matters because limits and usage are generally expressed in tokens, while users think in documents, sentences, or characters. Character counts are useful approximations for some local bounds, but they are not exact provider token counts.

For a token sequence \(x_1,\ldots,x_n\), the chain rule of probability gives:

\begin{equation}
p(x_1,\ldots,x_n)=\prod_{t=1}^{n}p(x_t\mid x_1,\ldots,x_{t-1}).
\end{equation}

An autoregressive language model approximates these conditional probabilities. During generation, previously supplied and generated tokens determine a distribution for the next token. This mathematical description explains why the complete input matters. Removing a policy passage changes the conditioning information; the model cannot be assumed to retain that passage from an earlier independent request unless the application or provider includes it again.

Prediction is not the same as factual verification. A likely continuation can contain a false date, an invented identifier, or a plausible but invalid tool argument. The application must decide which outputs can be accepted directly and which need an external check.

\section{A useful view of attention}
The Transformer introduced attention-based sequence processing \cite{transformer}. In a simplified attention operation, queries \(Q\), keys \(K\), and values \(V\) are matrices derived from token representations. One attention head computes:

\begin{equation}
\operatorname{Attention}(Q,K,V)=\operatorname{softmax}\left(\frac{QK^\top}{\sqrt{d_k}}\right)V.
\end{equation}

The dimension \(d_k\) is the width of the key vectors. A query–key dot product supplies a compatibility score. Softmax converts a row of scores into nonnegative weights summing to one, and those weights combine value vectors. This is a compact account of one operation, not a complete description of a modern provider model. Positional information, masking, multiple heads, feed-forward layers, and many architectural choices affect the full computation.

For the application designer, the relevant consequence is that context is processed through learned interactions. Including a sentence does not guarantee that the answer will use it correctly. Nor does an attention weight, by itself, establish a complete causal explanation for an answer. We therefore inspect task behavior through controlled inputs and external verification.

\section{Sampling and repeatability}
A model's raw output scores are often called logits. For positive temperature \(\tau\), a common sampling distribution is:

\begin{equation}
p_i=\frac{\exp(z_i/\tau)}{\sum_j\exp(z_j/\tau)}.
\end{equation}

Here \(z_i\) is the logit for candidate token \(i\). Lower temperature concentrates probability on higher-scoring candidates; higher temperature makes the distribution flatter. The mathematical expression applies for \(\tau>0\). An API's zero-temperature behavior is a separate implementation convention, often associated with greedy selection. A low temperature is not a universal guarantee of identical outputs across hardware, model revisions, or provider execution paths.

Top-\(p\) sampling retains a high-probability set whose cumulative mass reaches a threshold, then samples within that set. Provider support and parameter interactions vary. For an experiment, record the actual supported settings rather than copy a sampling recipe between models. Changing temperature, prompt, tools, and model at once makes it impossible to identify which change produced an observed difference.

\section{Prompting as interface design}
A prompt should establish the task, the available evidence, the required output, and what to do when evidence is insufficient. In the equipment service, “Be helpful” is under-specified. A more useful instruction says to select only catalog identifiers present in supplied records, state uncertainty about missing dates, and produce a reservation proposal without claiming that a booking has been committed.

Examples can clarify the intended mapping from inputs to outputs. A positive example alone may teach the model to always produce an item. Include a missing-evidence example if abstention is part of the contract. Keep examples consistent: if one example treats tomorrow as a fixed date and another uses the runtime clock, the application has left an important rule ambiguous.

Separate instructions from source content structurally. Use explicit fields for the user's request, retrieved records, and output requirements. This improves inspectability, but formatting is not an authorization boundary. Retrieved text remains untrusted even when enclosed in a field labeled “evidence.” Chapter 12 explains why tools require separate enforcement.

\section{Reasoning, explanations, and calculation}
Chain-of-thought prompting uses intermediate reasoning examples to elicit multi-step answers; Wei and colleagues studied its effects on arithmetic, commonsense, and symbolic tasks \cite{cot}. PAL instead uses model-generated programs with an interpreter for the computation \cite{pal}. These are different ways to allocate work. Neither lets the application accept an unchecked final claim merely because intermediate material looks plausible.

For an equipment fee, the model may extract “three days at 12 units per day” and request a calculator result. The calculator can establish that the product is 36. It cannot establish that the 12-unit rate is the applicable policy. Evidence selection and arithmetic verification are distinct responsibilities. A correct computation with a wrong input remains a wrong answer.

A verbal explanation is also an output to evaluate. If a model says it chose a camera because of battery life, an experiment might remove the battery information while holding other evidence constant. A changed answer supports sensitivity to that intervention; an unchanged answer does not prove the model ignored battery life in every context. CHIVE studies counterfactual prompt changes as a way to evaluate explanations of behavior \cite{chive}. Our classroom use is limited to input/output experiments and does not claim access to the model's internal mechanism.

\section{Worked example: three different validation failures}
Consider a required result with fields \texttt{item\_id}, \texttt{available}, and \texttt{evidence\_ids}. A response of \texttt{available: "yes"} violates a schema requiring a Boolean. A response with a valid Boolean and \texttt{item\_id: "C999"} may satisfy the schema but fail catalog membership. A response containing a real identifier and a real evidence ID may still misinterpret a policy that applies only to staff.

The checks form a sequence. Parsing determines whether there is a syntactically usable object. Schema validation checks field names, types, and local constraints. Domain validation checks membership and policy rules. Evidence review asks whether the cited material supports the particular claim. These checks need different error messages and different repair strategies.

If parsing fails, asking for a correctly formatted object may be sensible. If evidence is missing, repeating the same formatting instruction will not help. The system should retrieve missing evidence, ask a question, or abstain. Distinguishing failure types prevents a generic retry loop from repeatedly solving the wrong problem.

\section{Context, retrieval, and training}
Prompting changes the information and instructions supplied at inference time. Retrieval supplies external records relevant to the current request. Fine-tuning updates model parameters using training examples. These interventions operate at different points in the system and can be combined.

For frequently changing equipment availability, retrieval from the inventory service is the natural source of truth. Training a model to remember yesterday's inventory does not maintain today's bookings. For a stable output convention, examples or fine-tuning may improve adherence, but local validation remains necessary. Choosing among these approaches requires identifying whether the failure comes from missing knowledge, task interpretation, format adherence, or an unreliable downstream operation.

Structured output is helpful because software consumes objects more reliably than free prose. The supported schema features and behavior remain model- and provider-dependent; the accompanying labs use local validation even when requesting structured output. Consult the dated Gemini function-calling and compatibility references for the transport-specific details \cite{gemini-tools,gemini-openai}.

\section{Exercises}
\begin{enumerate}
\item Explain the difference between token probability and probability that a factual claim is correct. Give an example where a familiar phrase can be a false continuation.
\item For logits \((0,\ln 3)\) at temperature 1, compute the two softmax probabilities. Repeat at temperature 2 and explain the change.
\item A response cites a real policy ID but applies a staff rule to a student. Which checks pass and which fail? Propose the next action.
\item Write two prompt examples for equipment selection: one successful case and one insufficient-evidence case. State which behavior each example teaches.
\item Design a paired prompt experiment testing whether a misleading sentence changes an answer. Identify the outcome measure and two limits of the conclusion.
\end{enumerate}

\section{Further study and laboratory connection}
Read \cite{transformer} for the attention architecture and \cite{cot} for a historical reasoning intervention. Use \cite{chive} to distinguish explanatory hypotheses from behavioral tests. Lab 5 makes the model/calculator boundary explicit, and Lab 6 tests structured answers and evidence membership. Keep arithmetic correctness, source correctness, and semantic support separate in the lab report.

